\documentclass[10pt,a4paper]{amsart}
\usepackage[margin=19mm]{geometry}
\usepackage{amsmath,amssymb,mathtools}
\usepackage[scr=rsfs,cal=euler]{mathalfa}
\usepackage{xfrac}
\usepackage[unicode,hidelinks,bookmarks=false]{hyperref}
\newcommand{\ie}{i.e.\ }
\newcommand{\cf}{cf.\ }
\newcommand{\resp}{respectively\ }
\newcommand{\Subsection}{\subsection}
\providecommand{\nobreakdash}{\nobreak\hbox{-}}
\setlength{\emergencystretch}{2em}
\multlinegap0pt
\usepackage{amssymb}
\usepackage{xcolor}
\usepackage{algorithm}\usepackage{algpseudocode}
\usepackage[shortlabels]{enumitem}
\usepackage{float}
\newtheorem{lemma}{Lemma}
\newcommand{\E}{\mathbb{E}}
\title{\huge Geometric Milstein Scheme for Stochastic Differential Equations on SO(n) and SE(n)}
\author{Xi Wang, Victor Solo}
\date{Source: arXiv:2605.04480v1 (CC BY 4.0)}
\begin{document}
\maketitle

\noindent\emph{Source and licence.} \huge Geometric Milstein Scheme for Stochastic Differential Equations on SO(n) and SE(n). Version arXiv:2605.04480v1 (CC BY 4.0), distributed by arXiv under the Creative Commons Attribution 4.0 International licence, http://creativecommons.org/licenses/by/4.0/, as recorded in the arXiv licence metadata for this exact version and shown on the abstract page https://arxiv.org/abs/2605.04480v1. Full licence notice: \texttt{licenses/arxiv-2605.04480-CC-BY-4.0.txt}.\par\smallskip

\noindent\textit{Editorial scope.} Counted unit: the article's lemma lemma:error-sum-part (source0.tex lines 1188--1193). The article's complete original proof of it is delivered with it (source0.tex lines 1195--1220, one \texttt{proof} environment of 26 lines). No mathematics is written, repaired or re-derived: the statement and the proof are the article's own lines, in the article's own notation, and no retained line is substituted or re-flowed.  No further source-local context is needed: every cross-reference of every retained line resolves inside the two delivered spans.  Nothing was omitted from any retained span, so the package carries no omission record.  No retained line contains a citation, so the wrapper carries no bibliography and no bibliography entry.  No retained line contains a figure environment, an image-inclusion command or drawing code, so no image is delivered and the package declares no asset dependency.  Editorial formatting only, in addition: an AMS \texttt{amsart} preamble that restates the article's own declarations and macros used by the retained lines, namely the environments \texttt{lemma} on separate counters, together with the standard packages those lines need.  The retained environments print in this document as Lemma 1 in order of appearance, whereas the article prints the same items with the numbering of its own full document.  The complete native source of the article as distributed in the arXiv e-print for this exact version is bundled unchanged as \texttt{sources/199865/source0.tex}.
\begin{lemma}[Moment bound for partial‐sum]\label{lemma:error-sum-part} Let $\{E_k\}_{k=1}^{M}\subset\mathbb{R}^{d}$ be random vectors where $M = \mathcal{O}(\delta^{-1})$ for some time–step parameter $\delta > 0$. If $\E[\|E_m\|^{2}] = \mathcal{O}(\delta^{4})$ for every $m\le M$, then
\[
    \E\Bigl[\sup_{0\le m\le M} \|S_m\|^{2}\Bigr] 
    = \mathcal{O}(\delta^{2}).
\]
\end{lemma}

\begin{proof}
Assume that $\E[\|E_k\|^{2}] = \mathcal O(\delta^{4})$ for all $k$.
By Cauchy–Schwarz,
\[
    \|S_m\|^{2}
    =
    \Bigl\|\sum_{k=1}^{m}E_k\Bigr\|^{2}
    \le
    m\sum_{k=1}^{m}\|E_k\|^{2}
    \le
    M\sum_{k=1}^{M}\|E_k\|^{2},
\]
and therefore
\[
    \E\Bigl[\sup_{0\le m\le M}\|S_m\|^{2}\Bigr]
    \le
    M\sum_{k=1}^{M}\E[\|E_k\|^{2}]
    =
    M^{2}\mathcal O(\delta^{4})
    =
    \mathcal O(\delta^{2}).
\]
In either case we obtain  
$\E[\sup_{0\le m\le M}\|S_m\|^{2}] = \mathcal O(\delta^{2})$,  
which completes the proof.
\end{proof}

\end{document}
